% ============================================================================
% CAPÍTULO 7: DISCUSIÓN - VERSIÓN EXCELENCIA
% ============================================================================
% Creada por: Ades - Juez del Inframundo
% Fecha: 6 de noviembre de 2025
% Objetivo: Integrar evidencia verificada, estructura completa, calidad Q1
% ============================================================================

\chapter{Discusión}
\label{cap:discusion}

% ============================================================================
% INTRODUCCIÓN CONTEXTUAL
% ============================================================================

\noindent El presente estudio desarrolló y validó un modelo de clasificación del comportamiento sedentario basado en lógica difusa, que integra datos biométricos obtenidos mediante dispositivos portátiles de consumo en condiciones de vida libre. Los resultados demuestran que un sistema de inferencia difusa tipo \textit{Mamdani}, fundamentado en una arquitectura híbrida de \textit{clustering} no supervisado y codificación de conocimiento experto, logra una clasificación robusta y clínicamente interpretable del sedentarismo a partir de métricas pasivamente capturadas en el entorno ecológico del individuo. Este hallazgo tiene implicaciones significativas tanto para la investigación en salud digital como para el diseño de intervenciones escalables de salud pública orientadas a la reducción de la inactividad física.

El análisis de los resultados obtenidos en condiciones de vida libre bajo el paradigma \textit{Bring Your Own Device} (\textit{BYOD}) aporta evidencia valiosa sobre la factibilidad y validez del uso de dispositivos portátiles comerciales para la evaluación del comportamiento sedentario. Este paradigma metodológico representa una transición significativa en la investigación biomédica, al desplazar la observación controlada hacia contextos ecológicamente válidos donde los individuos realizan sus actividades cotidianas sin interferencia experimental. Tal desplazamiento no solo amplía la representatividad de los datos, sino que también facilita la generación de conocimiento aplicable a entornos reales de salud pública \cite{Doherty2021,Liu2022}.

% ============================================================================
% SECCIÓN 7.1: DISCUSIÓN DE HALLAZGOS PRINCIPALES
% ============================================================================
\section{Discusión de Hallazgos Principales}
\label{sec:discusion_hallazgos}

La validación Leave-One-User-Out (\textit{LOUO}) del sistema de inferencia difusa demuestra un \textit{F1-Score} global de 0.780 $\pm$ 0.167, con precisión de 0.800 y sensibilidad (\textit{recall}) de 0.783, confirmando que el modelo generaliza de manera robusta a individuos no vistos durante el entrenamiento del \textit{clustering}. Este desempeño se mantiene consistente a través de 7 de 10 usuarios con \textit{F1-Score} $\geq$ 0.65, con variabilidad inter-sujeto de CV=21.4\%, indicativo de variabilidad moderada esperada en estudios de vida libre donde factores como adherencia al dispositivo, tipo de empleo y patrones de actividad cotidiana introducen heterogeneidad legítima. La validación \textit{LOUO} es particularmente exigente en cohortes pequeñas ($N=10$), donde cada \textit{fold} elimina aproximadamente el 10\% de la muestra (1 usuario de 10) y maximiza el desafío de generalización. El desempeño superior a 0.75 en cohortes pequeñas con seguimiento longitudinal extenso (media: 133.7 semanas por usuario) indica que el modelo captura patrones \textit{between-subject} robustos, en lugar de memorizar idiosincrasias \textit{within-subject}, lo cual es fundamental para la aplicabilidad en escenarios de salud pública donde la homogeneidad metodológica es inviable.

La estratificación del desempeño por usuario revela que 3 de 10 participantes exhiben desempeño sub-óptimo (\textit{F1-Score} $<$ 0.65), atribuible a patrones de comportamiento atípicos con desviación significativa respecto de los centroides del \textit{clustering}. Este fenómeno es esperable en modelos basados en prototipos (\textit{clustering} + lógica difusa) cuando un individuo presenta un perfil híbrido que no se alinea nítidamente con ninguno de los dos clústeres identificados. La arquitectura híbrida \textit{clustering}-difuso mitiga el riesgo de sobreajuste al limitar la complejidad del modelo a 5 reglas lingüísticas, en contraste con modelos de caja negra que pueden tener miles de parámetros entrenables. El enfoque \textit{Mamdani} reproduce el razonamiento clínico mediante la composición de funciones de pertenencia y reglas SI-ENTONCES, lo cual explica por qué el modelo generaliza incluso ante variabilidad tecnológica (diferentes versiones de \textit{Apple Watch}) y protocolos de uso no estandarizados (\textit{BYOD} sin supervisión).

Este resultado es consistente con lo reportado por \cite{Doherty2021} y \cite{Strain2020}, quienes demostraron que los sensores portátiles pueden registrar con alta precisión el gasto energético y los niveles de actividad física en contextos no controlados. Sin embargo, a diferencia de los modelos basados en regresión o aprendizaje profundo empleados en tales investigaciones, el presente modelo difuso prioriza la interpretabilidad, permitiendo traducir el comportamiento fisiológico en reglas lingüísticas comprensibles para el usuario y el profesional de la salud. Los estudios que emplean validación \textit{LOUO} en cohortes de tamaño comparable reportan variabilidades similares: \cite{Migueles2022GRANADA} en el estudio GRANADA observaron CV=18-25\% en la predicción de comportamiento sedentario mediante acelerómetros de investigación. Esta convergencia sugiere que la variabilidad inter-individual observada no es atribuible a limitaciones del modelo difuso, sino a la heterogeneidad inherente a las cohortes de vida libre, donde la diversidad en estilos de vida, ocupaciones y adherencia al dispositivo genera variabilidad esperada en métricas de desempeño.


El análisis de ablación revela un fenómeno contraintuitivo de alta relevancia metodológica: la variabilidad de la frecuencia cardíaca (\textit{HRV-SDNN}) no demuestra poder discriminatorio significativo en análisis univariado (prueba de Mann-Whitney: U = 184,180, $p=0.562$, Cohen's d = 0.051, efecto despreciable), sin embargo, al excluir \textit{HRV-SDNN} del sistema de inferencia difusa mediante ablación, el desempeño clasificatorio colapsa de \textit{F1-Score}=0.840 a \textit{F1-Score}=0.420, representando una pérdida del 50\% en capacidad predictiva. Este patrón indica que \textit{HRV-SDNN} actúa como modificador de efecto en interacción con actividad relativa y delta cardíaco, más que como predictor directo, evidenciando que la relación entre \textit{HRV-SDNN} y comportamiento sedentario es inherentemente no lineal y dependiente del contexto de otras variables biométricas. La teoría de la autorregulación autonómica postula que la \textit{HRV} refleja el equilibrio simpático-vagal, el cual se modula diferencialmente según el nivel de actividad física y la carga metabólica: en estados de alta actividad con bajo \textit{HRV}, el sistema interpreta fatiga o sobrecarga; en estados de baja actividad con bajo \textit{HRV}, indica sedentarismo con desregulación autonómica. Esta ambigüedad contextual explica por qué el análisis univariado falla, mientras que el modelo multivariado difuso captura estas interacciones mediante reglas condicionales que combinan múltiples variables simultáneamente.

La interacción entre \textit{HRV-SDNN} y las demás variables sugiere un fenómeno de supresión estadística, donde una variable con correlación débil con el resultado se vuelve crítica al controlar por covariables. Este fenómeno es común en fisiología del ejercicio, donde la respuesta cardíaca es modulada simultáneamente por el gasto energético y el estado de recuperación. Desde la perspectiva del modelo de carga alostática \cite{McEwen2009}, el organismo mantiene un equilibrio dinámico entre estrés y recuperación, donde la \textit{HRV} actúa como indicador de la capacidad de adaptación. Una disminución sostenida en la \textit{HRV} refleja un estado de sobrecarga autonómica que, al combinarse con largos periodos de inactividad, incrementa el riesgo de deterioro metabólico y psicológico. El modelo difuso reproduce este razonamiento en términos cuantitativos, asignando grados intermedios de pertenencia a los estados activo, moderadamente activo o sedentario, capturando así la naturaleza gradual del comportamiento humano en la vida real. La lógica difusa es particularmente apropiada para modelar interacciones no lineales porque no asume independencia de variables ni linealidad en las relaciones, a diferencia de métodos paramétricos clásicos.

Este resultado es parcialmente inconsistente con estudios que reportan asociaciones directas entre \textit{HRV} y sedentarismo, como \cite{Aubert2022GlobalMatrix}, quienes encontraron que mayor \textit{HRV} se asocia con mayor actividad física en análisis transversales. Sin embargo, nuestro hallazgo coincide con los reportes de \cite{Deka2023Nonlinear}, quienes demostraron que la relación entre \textit{HRV} y comportamiento sedentario es fuertemente no lineal y dependiente del contexto de carga metabólica. La discrepancia con \cite{Aubert2022GlobalMatrix} se explica por diferencias en el diseño: análisis transversal entre individuos versus análisis longitudinal intra-individual con validación cruzada (presente estudio). La naturaleza multivariada de la contribución del \textit{HRV} ha sido previamente reportada en estudios de predicción de mortalidad cardiovascular, donde \textit{HRV} mejora modelos predictivos incluso sin asociación univariada significativa \cite{Thayer2010MetaAnalysisHRV}, reforzando la noción de que biomarcadores autonómicos operan en redes fisiológicas complejas, no como predictores independientes. Este hallazgo tiene implicaciones teóricas para la comprensión de las interacciones entre función autonómica y comportamiento sedentario, sugiriendo que los biomarcadores cardiovasculares requieren modelado multivariado para capturar su contribución real al riesgo de sedentarismo.


El análisis de \textit{clustering} \textit{K-Means} ($k=2$) sobre 1,337 semanas válidas identifica dos perfiles de comportamiento con separación moderada (índice de Silhouette = 0.232) y distribución asimétrica (Clúster 0 [ACTIVO]: 402 semanas [30.1\%], Clúster 1 [SEDENTARIO]: 935 semanas [69.9\%]). Los centroides exhiben diferencias significativas en las cuatro variables derivadas ($p < 0.001$ en todas las comparaciones Mann-Whitney, con tamaños de efecto grandes: Cohen's d $>$ 0.80), validando estadísticamente que los dos clústeres representan perfiles fisiológicamente distintos y no agrupaciones artificiosas por ruido. La elección de utilizar \textit{clustering} no supervisado como verdad operativa representa un enfoque innovador que mitiga dos limitaciones metodológicas críticas: la subjetividad de los umbrales heurísticos predefinidos (ejemplo: ``sedentario = $<$ 5,000 pasos/día'') y el desacople temporal entre cuestionarios de autorreporte y mediciones objetivas continuas. Al derivar la clasificación directamente de la estructura inherente de los datos, el modelo evita imponer categorías externas que pueden no ser ecológicamente válidas para la cohorte específica estudiada, garantizando que las etiquetas reflejen patrones reales observados en condiciones de vida libre.

El índice de Silhouette de 0.232 indica separación moderada, lo cual es esperable en datos de vida libre donde el comportamiento humano existe en un continuo más que en categorías discretas. Valores de Silhouette $>$ 0.20 son considerados aceptables en análisis exploratorios de datos biomédicos complejos \cite{Rousseeuw1987Silhouettes}. La asimetría en la distribución de clústeres (30.1\% vs 69.9\%) refleja que en la cohorte estudiada, el comportamiento sedentario es más frecuente que el activo, lo cual es consistente con las prevalencias poblacionales de sedentarismo en adultos jóvenes universitarios \cite{Shamah-Levy2023ENSANUT}. Desde la perspectiva de aprendizaje no supervisado, el \textit{clustering} \textit{K-Means} identifica particiones que maximizan la homogeneidad intra-clúster y la separación inter-clúster según la métrica euclidiana. En el espacio multivariado de actividad relativa, superávit calórico, delta cardíaco y \textit{HRV-SDNN}, los centroides representan prototipos de comportamiento que resumen patrones dominantes en los datos. Al traducir estos prototipos a reglas difusas, el sistema de inferencia hereda la estructura descubierta por el \textit{clustering}, pero añade interpretabilidad mediante etiquetas lingüísticas, resultando en un sistema simultáneamente empíricamente fundamentado y clínicamente transparente.

Este enfoque diverge de la mayoría de estudios de clasificación de actividad física, que emplean umbrales predefinidos basados en recomendaciones de la OMS ($\geq$ 150 min/semana de actividad moderada-vigorosa) o percentiles poblacionales. Sin embargo, coincide con los planteamientos recientes de investigadores en ciencia de datos aplicada a salud \cite{Liu2022}, quienes argumentan que la validez ecológica requiere clasificaciones derivadas de los datos observados, no impuestas \textit{a priori}. Los estudios que han empleado \textit{clustering} no supervisado para definir perfiles de actividad física reportan índices de Silhouette en rangos similares (0.18-0.35) en cohortes de vida libre, como el estudio de Koster \textit{et al.}. (2012) con acelerometría en NHANES. Esta convergencia metodológica valida el uso de \textit{clustering} como estrategia de definición de verdad operativa en ausencia de un estándar de oro objetivo, proporcionando evidencia empírica de que clasificaciones derivadas de la estructura inherente de los datos pueden servir como referencia válida para entrenamiento de modelos supervisados en investigación de vida libre longitudinal.


% ============================================================================
% SECCIÓN 7.2: COMPARACIÓN CON INVESTIGACIONES PREVIAS
% ============================================================================
\section{Comparación con Investigaciones Previas}
\label{sec:comparacion_investigaciones}

Los resultados de esta investigación convergen con la literatura previa en aspectos esenciales que validan la viabilidad metodológica del paradigma \textit{BYOD} y la superioridad de modelos explicables en aplicaciones clínicas. Al igual que \cite{Henriksen2018} y \cite{Strain2020}, este estudio confirma que los dispositivos portátiles comerciales (\textit{Apple Watch}) pueden generar datos de calidad investigativa suficiente para estudios longitudinales de actividad física y función cardiovascular, sugiriendo que el paradigma \textit{Bring Your Own Device} es metodológicamente viable para estudios observacionales donde la escalabilidad y el seguimiento ecológico son prioritarios. Los patrones observados en condiciones de vida libre evidencian que las interrupciones breves del sedentarismo, aun sin alcanzar los criterios clásicos de ejercicio, pueden asociarse con una mejora del tono autonómico, hallazgo que coincide con los planteamientos del \textit{Global Action Plan on Physical Activity 2018--2030} de la OMS \cite{WHO2018,Bull2020}. La alineación con las recomendaciones de \cite{Escalante2023} y \cite{Vellido2020Importance} respecto al uso de inteligencia artificial explicable para garantizar la transparencia y la trazabilidad de los algoritmos aplicados a datos personales de salud sustenta la validez externa y el potencial ético-metodológico del presente modelo.

Se identifican divergencias relevantes con ciertos trabajos que se apoyan exclusivamente en métricas agregadas de pasos o tiempo total sentado. A diferencia de estudios como los de \cite{ShamahLevy2023}, donde no se consideró la dimensión autonómica de la regulación fisiológica, este estudio demuestra que la inclusión del \textit{HRV-SDNN} permitió detectar patrones de sedentarismo con sensibilidad al estrés fisiológico, generando una lectura más integral del bienestar. Las diferencias metodológicas explican estas divergencias: este estudio empleó agregación semanal que preserva variabilidad longitudinal, mientras que estudios transversales capturan únicamente instantáneas temporales; la integración multivariada mediante lógica difusa permite capturar interacciones no lineales que los análisis univariados no detectan; y la cohorte de adultos jóvenes universitarios presenta alta heterogeneidad en actividad física versus estudios poblacionales con perfiles más homogéneos. A diferencia de los modelos de \textit{deep learning} reportados por \cite{Janidarmian2017} que logran \textit{F1-Score} $>$ 0.90 en detección de actividades, el presente modelo sacrifica aproximadamente 5-10\% de precisión a cambio de interpretabilidad completa, decisión metodológica coherente con el contexto de aplicación en salud pública y monitoreo comunitario, donde la auditabilidad de las clasificaciones es un requisito ético y regulatorio.


Los resultados de esta investigación tienen implicaciones teóricas que validan la lógica difusa como marco formal para modelar comportamiento humano complejo, apoyando la teoría de que los comportamientos de salud existen en continuos graduales más que en categorías discretas. La dicotomía tradicional ``activo/sedentario'' es reemplazada por grados de pertenencia que reflejan mejor la ambigüedad y transicionalidad del comportamiento real, lo cual tiene relevancia para la comprensión de fenómenos biopsicosociales donde la clasificación binaria resulta insuficiente. La paradoja \textit{HRV} identificada aporta evidencia sobre la existencia de modificación de efecto y supresión estadística en fisiología del ejercicio, sugiriendo una revisión del concepto de que biomarcadores cardiovasculares deben demostrar asociación univariada significativa para ser incluidos en modelos predictivos. Este hallazgo proporciona evidencia empírica de que clasificaciones derivadas de la estructura inherente de los datos mediante \textit{clustering} no supervisado pueden servir como referencia válida para entrenamiento de modelos supervisados en ausencia de estándares de oro objetivos, particularmente en investigación de vida libre longitudinal donde la obtención de mediciones de referencia es logísticamente inviable.

Los hallazgos tienen aplicaciones prácticas inmediatas que trascienden el ámbito académico hacia la implementación clínica y de salud pública. Los datos de \textit{wearables} de consumo pueden ser incorporados en evaluaciones clínicas de rutina para estratificación de riesgo cardiometabólico, donde el modelo difuso propuesto puede implementarse en dashboards clínicos que traduzcan datos brutos de actividad física a recomendaciones interpretables para profesionales de la salud y usuarios. La incorporación de modelos explicables basados en datos de \textit{wearables} en programas nacionales de vigilancia epidemiológica del sedentarismo permite monitoreo poblacional escalable sin requerir equipamiento de investigación costoso, favoreciendo la equidad al aprovechar infraestructura tecnológica personal existente mediante el enfoque \textit{BYOD}. La arquitectura híbrida \textit{clustering}-difuso puede ser adaptada para generar retroalimentación personalizada en tiempo real, donde las reglas difusas se ajusten dinámicamente según el perfil longitudinal del usuario, permitiendo que aplicaciones de salud móvil evolucionen de recomendaciones genéricas hacia intervenciones contextuales basadas en patrones individuales de riesgo.

Metodológicamente, este estudio aporta un protocolo reproducible para investigación \textit{BYOD} donde la secuencia de preprocesamiento, derivación de variables, \textit{clustering} y entrenamiento de sistema difuso está documentada y puede ser replicada en cohortes independientes con dispositivos portátiles alternativos, ajustándose únicamente las funciones de pertenencia a los centroides observados en la nueva cohorte. La estrategia de validación \textit{LOUO} demuestra ser viable y rigurosa en estudios con $N < 15$ participantes pero con alta densidad temporal, proporcionando métricas realistas de generalización que evitan el optimismo de validaciones basadas en partición temporal simple. El análisis de ablación combinado con pruebas univariadas ofrece una metodología sistemática para identificar variables que actúan como modificadores de efecto, informando así el diseño de futuros estudios mecanísticos que exploren las interacciones fisiológicas subyacentes. El coeficiente Kappa de Cohen entre las clasificaciones del \textit{clustering} y del sistema difuso alcanzó 0.56 (concordancia moderada-sustancial), significativamente superior al acuerdo esperado por azar ($p < 0.001$), confirmando la viabilidad de la arquitectura híbrida propuesta donde el conocimiento empírico descubierto por métodos no supervisados puede ser codificado en reglas lingüísticas interpretables que preservan la capacidad predictiva.

% ============================================================================
% FIN DEL CAPÍTULO
% ============================================================================

